\documentclass[11pt]{article}
\usepackage[margin=1in]{geometry}
\usepackage{amsmath,amssymb,amsthm}
\usepackage{booktabs}
\usepackage[hidelinks]{hyperref}

\newtheorem{theorem}{Theorem}
\newtheorem{lemma}[theorem]{Lemma}

\newcommand{\supp}{\operatorname{supp}}
\newcommand{\size}{\operatorname{size}}
\newcommand{\F}{F}

\title{\textbf{The Provenance Forest}\\[1mm]
\large Poker's generalisation step without enumerating refutations}
\author{Oru\v{z}i \and Zarathustra Goertzel}
\date{7 November 2026}

\begin{document}
\maketitle

\begin{abstract}
Poker generalises an example by collecting, for every refutation of the example
within the clause limit, the set of metasubstitutions the refutation uses. Vanilla
enumerates these refutations one at a time, and their number grows exponentially
with the clause limit. The provenance forest computes the same sets by dynamic
programming: it solves each first-order goal once for each exact size, and builds
the sets of a refutation from those of its parts. On Poker's $a^nb^n$ experiment it
returns exactly Poker's sub-hypotheses in every case compared, for example the same
18 at $a^4b^4$ with clause limit 7 in 3.4\,ms where Poker takes 218\,s.
\end{abstract}

\section{The problem}

Poker's generalisation step takes an example $e$, background knowledge $B$, a set of
metarules, a signature (the target predicate and the invented predicates) and a
clause limit $K$. Each refutation of $e$ resolves goals either against $B$ or against
an instance of a metarule; each such instance is a \emph{metasubstitution} $\sigma$,
the metarule with its second-order variables bound. The \emph{support} of a refutation
is the set of metasubstitutions it uses. Poker's sub-hypotheses are the supports of
the refutations of $e$ with at most $K$ steps.

Vanilla finds them by SLD-resolution with the metasubstitutions found so far as an
argument of every call. Distinct refutations share most of their sub-refutations,
but an untabled prover derives each shared part again in every refutation that
contains it. Tabling the meta-interpreter does not recover the sharing: two calls on
the same goal with different stores are different calls. On $a^4b^4$ at clause limit
7 Poker spends 5.13 billion inferences to produce 18 sets.

\section{The decomposition}

A refutation is a tree. Its root is a metasubstitution $\sigma$ whose head matches
the goal; its children are refutations of the body literals of $\sigma$; its leaves
are background facts. Size and support are both compositional:
\[
  \size(t) = 1 + \sum_i \size(t_i), \qquad
  \supp(t) = \{\sigma\} \cup \bigcup_i \supp(t_i),
\]
with size $0$ and support $\emptyset$ at a background leaf. The size counts steps
with multiplicity, while the support is a set, as in Poker's store. Neither depends
on the rest of the tree, so the family of supports of a goal's refutations of size
exactly $n$,
\[
  \F_n(g) = \{\, \supp(t) : t \text{ a refutation of } g,\ \size(t) = n \,\},
\]
satisfies
\begin{align*}
  \F_0(g)     &= \{\emptyset\} \text{ if } g \in B, \text{ and } \emptyset \text{ otherwise},\\
  \F_{n+1}(g) &= \bigcup_{\sigma : g \leftarrow g_1, \ldots, g_m}\
                 \bigcup_{n_1 + \cdots + n_m = n}
                 \bigl\{\, \{\sigma\} \cup S_1 \cup \cdots \cup S_m :
                          S_i \in \F_{n_i}(g_i) \,\bigr\},
\end{align*}
and the generalisation of $e$ at clause limit $K$ is $\bigcup_{n=1}^{K} \F_n(e)$.
Each $\F_n(g)$ is computed once and shared by every refutation that contains $g$ at
size $n$.

\begin{theorem}[the recurrence characterises the supports]
For a ground derivation forest (goals, background goals, and edges
$g \leftarrow \sigma, [g_1, \ldots, g_m]$), the family $\F$ of supports of derivation
trees satisfies the recurrences above, and every family satisfying them equals $\F$.
\end{theorem}

\begin{proof}[Proof sketch]
That $\F$ satisfies the recurrences is the decomposition of a tree at its root. For
uniqueness, let $G$ satisfy them and show $G_n(g) = \F_n(g)$ for all $g$ by strong
induction on $n$: at $n = 0$ both are given by $B$; at $n + 1$ both sides apply the
same recurrence to families of sizes $n_i \le n$, which agree by induction.
\end{proof}

\begin{lemma}[pruning loses nothing]
In the recurrence, a splitting $n_1 + \cdots + n_m = n$ contributes a set only if
each body literal $g_i$ has a derivation of size $n_i$. Hence restricting the
splittings to those in which every literal is known to be derivable at its share
keeps every set of $\F_{n+1}(g)$.
\end{lemma}

Both statements are machine-checked in Lean~4 with Mathlib, using no axioms beyond
Lean's standard ones, together with positive and negative examples.

\section{From ground forests to MIL goals}

In meta-interpretive learning the forest is not given in advance. A goal is
$P(x, Y)$ with $x$ known; the body of Chain, $Q(X, Z), R(Z, Y)$, introduces an
existential $Z$ that is known only once $Q(x, Z)$ is solved. The forest therefore
works in two passes over memo tables keyed by goal and size.

\paragraph{Pass 1: endpoints.} $\mathit{reach}(P, x, n)$ is the set of endpoints $y$
such that $P(x, y)$ has a refutation of size exactly $n$. A body is solved left to
right: the size $n - 1$ is split among the literals, and the endpoints of each
literal at its share instantiate the next literal's input. A literal of size $0$ is
answered by $B$; a literal of positive size is a target or invented predicate,
answered by $\mathit{reach}$ itself.

\paragraph{Pass 2: supports.} $\mathit{fam}(P, x, y, n)$ is $\F_n(P(x,y))$, built by
the recurrence, but only over the splittings pass 1 found feasible. By the pruning
lemma this loses nothing.

\paragraph{Invented predicates.} Each memo key also carries which invented predicates
are available: the $i$-th invented predicate may head a clause only once the
$(i{-}1)$-th is available. This names inventions canonically, so that supports
differing only by a renaming of invented predicates are not produced twice. A goal
of the ground forest is then a goal together with this state, and the theorem
applies unchanged.

\paragraph{Metarule constraints.} Poker's constraints, such as $P \neq Q$ in Chain,
are tests on a single metasubstitution. They are applied as soon as a
metasubstitution is built, comparing symbols with \texttt{==} so that a symbol not
yet bound passes. A constraint on one metasubstitution removes edges from the forest
and does not interact with the decomposition.

\section{Cost}

The memo tables hold one entry per goal, size and invention state. Pass 1 costs,
per entry, one step for each metarule, choice of symbol, splitting of the size and
intermediate endpoint. Pass 2 builds only sets that are parts of sub-hypotheses:
every set it builds is contained in one of the output's sets. Poker's cost grows with
the number of refutations instead, and of the partial ones that fail, a number that
grows exponentially with $K$ even when the number of distinct supports does not: at
$a^4b^4$, limit 7, there are 18 supports.

Table~\ref{tab:results} reports Poker's \texttt{generalise/3} and the forest on the
same problems, in SWI-Prolog 10.1.9. The sets of sub-hypotheses are identical in
every row. At limits too small to fit any refutation both return none (Poker's
\texttt{generalise/3} fails, which the comparison reads as none).

\begin{table}[h]
\centering
\begin{tabular}{lrrrrr}
\toprule
Task & Sub-hypotheses & \multicolumn{2}{c}{Poker} & \multicolumn{2}{c}{Forest}\\
     &               & time & inferences & time & inferences\\
\midrule
$a^1b^1$, $K=1$ & 1     & 0.9\,ms  & 5,956          & 0.2\,ms  & 1,804\\
$a^2b^2$, $K=3$ & 4     & 9.3\,ms  & 199,702        & 0.4\,ms  & 6,399\\
$a^2b^2$, $K=5$ & 1,449 & 1.68\,s  & 35,883,880     & 6.9\,ms  & 79,931\\
$a^3b^3$, $K=5$ & 10    & 1.20\,s  & 29,253,930     & 1.1\,ms  & 23,176\\
$a^3b^3$, $K=7$ & 9,941 & 219\,s   & 5,348,394,103  & 124\,ms  & 711,019\\
$a^4b^4$, $K=7$ & 18    & 218\,s   & 5,134,352,712  & 3.4\,ms  & 65,625\\
\bottomrule
\end{tabular}
\caption{Poker against the forest on Poker's $a^nb^n$ experiment (single runs; the
inference counts are exact).}
\label{tab:results}
\end{table}

On its own the forest takes 7.1\,ms on $a^5b^5$ at $K = 9$ and 32\,ms on $a^7b^7$ at
$K = 13$.

\section{Scope}

\begin{itemize}
\item The decomposition assumes that a refutation resolves goals against background
  knowledge and metarules only, not against clauses of the hypothesis being built,
  so that every step instantiates a metarule and the clause limit bounds the steps.
  This is the setting of the $a^nb^n$ experiment: there Poker finds no refutation of
  $a^3b^3$ at limit 4, although a grammar of three clauses derives it in five steps.
  When hypothesis clauses can be reused, size is no longer additive; extending the
  forest to that setting is open.
\item The forest replaces the generalisation step only. Poker's other stages are
  unchanged.
\item Tasks are currently given in the forest's own encoding; reading Vanilla
  experiment files directly is the next step.
\end{itemize}

\section{Running it}

The forest is in \texttt{lib/forest}; tasks and scripts are in
\texttt{data/forest\_examples}. From Vanilla's root, with Poker selected as the
learner in \texttt{configuration.pl},
\begin{center}
\texttt{swipl data/forest\_examples/poker\_comparison.pl N K}
\end{center}
runs both learners on $a^Nb^N$ at clause limit $K$ and compares their sub-hypotheses
as sets of sets. The forest was first written in MeTTa, where each memo table is a
space of shared immutable terms indexed by their arguments; the files here are its
port to SWI-Prolog, with the memo tables as asserted facts keyed by term hash or as
tabled predicates.

\end{document}
